% GENERATED FILE. Edit canonical lessons, then run npm run generate:books.
% canonical-source-hash: fnv1a64:027b2b5345186a9a
% canonical-lessons: MR-C535-citrakala, MR-C535-kalakar, MR-C535-silpa, MR-C535-kuncla, MR-C535-chata, MR-R535-review-528-535

\chapter{Painting and Sculpture}
\label{ch:mr-painting-and-sculpture}
\hlchaptermodality{535}

\begin{chapteropening}
\textbf{By the end of this chapter:} \emph{I can talk about painting, sculpture and artists.}

Complete the last lesson of chapter 535: I can talk about painting, sculpture and artists.
\end{chapteropening}

\section[\texorpdfstring{\mr{चित्रकला} (citrakalā)}{चित्रकला (citrakalā)}]{\mr{चित्रकला} (citrakalā) — painting as an art}

\label{lesson:MR-C535-citrakala}



\begin{quote}
Before the new one: say the Marathi for a tent, then the Marathi for a campfire.
\end{quote}

\subsection*{You'll want to know: \mr{चित्रकला}}
\textbf{\mr{चित्रकला}} — \emph{citrakalā} — \textquotedblleft{}painting as an art\textquotedblright{}.

\begin{quote}
\mr{मीराला} \mr{लहानपणापासून} \mr{चित्रकलेची} \mr{आवड} \mr{आहे}.
\end{quote}

\begin{quote}
\emph{mīrālā lahānpaṇāpāsūn citrakalecī āvaḍ āhe.}
\end{quote}

\begin{quote}
Meera has loved painting since childhood.
\end{quote}

Feminine.

\begin{cousinweb}
\textbf{\mr{चित्रकला}} (\emph{citrakalā}) is a tatsama compound of Sanskrit \emph{citra}, \textquotedblleft{}picture\textquotedblright{}, and \emph{kalā}, \textquotedblleft{}art\textquotedblright{} (Monier-Williams).
\end{cousinweb}

\begin{grammarlens}[title={What you've built}]
One more word for describing what you lived and felt.
\end{grammarlens}

\subsection*{Your turn}
\begin{itemize}
\raggedright
  \item \emph{Say it:} \emph{citrakalā}
  \item \emph{Say it:} \emph{citrakalā}, once more
  \item \emph{Say it:} \emph{mīrālā lahānpaṇāpāsūn citrakalecī āvaḍ āhe.} — Meera has loved painting since childhood.
  \item \emph{Recall:} say the Marathi for a tent, then the Marathi for a campfire, then say \emph{citrakalā} again
\end{itemize}

\begin{tcolorbox}[breakable,colback=teal!4,colframe=teal!35!black,title={Before you move on}]
What does \mr{चित्रकला} mean? (\textquotedblleft{}Painting as an art\textquotedblright{}.) Say it once more.
\end{tcolorbox}
\section[\texorpdfstring{\mr{कलाकार} (kalākār)}{कलाकार (kalākār)}]{\mr{कलाकार} (kalākār) — an artist}

\label{lesson:MR-C535-kalakar}



\begin{quote}
Before the new one: say the Marathi for a campfire, then the Marathi for painting as an art.
\end{quote}

\subsection*{You'll want to know: \mr{कलाकार}}
\textbf{\mr{कलाकार}} — \emph{kalākār} — \textquotedblleft{}an artist\textquotedblright{}.

\begin{quote}
\mr{हा} \mr{कलाकार} \mr{मातीपासून} \mr{सुंदर} \mr{वस्तू} \mr{बनवतो}.
\end{quote}

\begin{quote}
\emph{hā kalākār mātīpāsūn sundar vastū banavto.}
\end{quote}

\begin{quote}
This artist makes beautiful things out of clay.
\end{quote}

Masculine.

One word for a man or a woman.

\begin{cousinweb}
\textbf{\mr{कलाकार}} (\emph{kalākār}) is a compound built from Sanskrit \emph{kalā}, \textquotedblleft{}art\textquotedblright{}, and \emph{kāra}, \textquotedblleft{}maker\textquotedblright{} (Monier-Williams).
\end{cousinweb}

\begin{grammarlens}[title={What you've built}]
One more word for describing what you lived and felt.
\end{grammarlens}

\subsection*{Your turn}
\begin{itemize}
\raggedright
  \item \emph{Say it:} \emph{kalākār}
  \item \emph{Say it:} \emph{kalākār}, once more
  \item \emph{Say it:} \emph{hā kalākār mātīpāsūn sundar vastū banavto.} — This artist makes beautiful things out of clay.
  \item \emph{Recall:} say the Marathi for a campfire, then the Marathi for painting as an art, then say \emph{kalākār} again
\end{itemize}

\begin{tcolorbox}[breakable,colback=teal!4,colframe=teal!35!black,title={Before you move on}]
What does \mr{कलाकार} mean? (\textquotedblleft{}An artist\textquotedblright{}.) Say it once more.
\end{tcolorbox}
\section[\texorpdfstring{\mr{शिल्प} (śilpa)}{शिल्प (śilpa)}]{\mr{शिल्प} (śilpa) — a sculpture}

\label{lesson:MR-C535-silpa}



\begin{quote}
Before the new one: say the Marathi for painting as an art, then the Marathi for an artist.
\end{quote}

\subsection*{You'll want to know: \mr{शिल्प}}
\textbf{\mr{शिल्प}} — \emph{śilpa} — \textquotedblleft{}a sculpture\textquotedblright{}.

\begin{quote}
\mr{बागेत} \mr{दगडाचं} \mr{एक} \mr{मोठं} \mr{शिल्प} \mr{आहे}.
\end{quote}

\begin{quote}
\emph{bāget dagḍācaṁ ek moṭhaṁ śilpa āhe.}
\end{quote}

\begin{quote}
There is a big stone sculpture in the garden.
\end{quote}

Neuter.

\begin{cousinweb}
\textbf{\mr{शिल्प}} (\emph{śilpa}) is a tatsama, taken from Sanskrit \emph{śilpa}, \textquotedblleft{}artistic work, a craft\textquotedblright{} (Monier-Williams).
\end{cousinweb}

\begin{grammarlens}[title={What you've built}]
One more word for describing what you lived and felt.
\end{grammarlens}

\subsection*{Your turn}
\begin{itemize}
\raggedright
  \item \emph{Say it:} \emph{śilpa}
  \item \emph{Say it:} \emph{śilpa}, once more
  \item \emph{Say it:} \emph{bāget dagḍācaṁ ek moṭhaṁ śilpa āhe.} — There is a big stone sculpture in the garden.
  \item \emph{Recall:} say the Marathi for painting as an art, then the Marathi for an artist, then say \emph{śilpa} again
\end{itemize}

\begin{tcolorbox}[breakable,colback=teal!4,colframe=teal!35!black,title={Before you move on}]
What does \mr{शिल्प} mean? (\textquotedblleft{}A sculpture\textquotedblright{}.) Say it once more.
\end{tcolorbox}
\section[\texorpdfstring{\mr{कुंचला} (kuñclā)}{कुंचला (kuñclā)}]{\mr{कुंचला} (kuñclā) — a paintbrush}

\label{lesson:MR-C535-kuncla}



\begin{quote}
Before the new one: say the Marathi for an artist, then the Marathi for a sculpture.
\end{quote}

\subsection*{You'll want to know: \mr{कुंचला}}
\textbf{\mr{कुंचला}} — \emph{kuñclā} — \textquotedblleft{}a paintbrush\textquotedblright{}.

\begin{quote}
\mr{कलाकाराने} \mr{बारीक} \mr{कुंचल्याने} \mr{डोळे} \mr{रंगवले}.
\end{quote}

\begin{quote}
\emph{kalākārāne bārīk kuñclyāne ḍoḷe raṅgavle.}
\end{quote}

\begin{quote}
The artist painted the eyes with a fine brush.
\end{quote}

Masculine.

The oblique is \mr{कुंचल्या}-: \mr{कुंचल्याने}.

\begin{cousinweb}
\textbf{\mr{कुंचला}} (\emph{kuñclā}) probably belongs with Marathi \emph{kũcā} and \emph{kũcī}, \textquotedblleft{}brush\textquotedblright{}, which CDIAL sets under Sanskrit \emph{kūrca}, \textquotedblleft{}a bunch of grass\textquotedblright{}, whose derivatives meant a painter's brush (CDIAL 3408). CDIAL lists no \mr{कुंचला} itself, so the -\mr{ला} form is not settled.
\end{cousinweb}

\begin{grammarlens}[title={What you've built}]
One more word for describing what you lived and felt.
\end{grammarlens}

\subsection*{Your turn}
\begin{itemize}
\raggedright
  \item \emph{Say it:} \emph{kuñclā}
  \item \emph{Say it:} \emph{kuñclā}, once more
  \item \emph{Say it:} \emph{kalākārāne bārīk kuñclyāne ḍoḷe raṅgavle.} — The artist painted the eyes with a fine brush.
  \item \emph{Recall:} say the Marathi for an artist, then the Marathi for a sculpture, then say \emph{kuñclā} again
\end{itemize}

\begin{tcolorbox}[breakable,colback=teal!4,colframe=teal!35!black,title={Before you move on}]
What does \mr{कुंचला} mean? (\textquotedblleft{}A paintbrush\textquotedblright{}.) Say it once more.
\end{tcolorbox}
\section[\texorpdfstring{\mr{छटा} (chaṭā)}{छटा (chaṭā)}]{\mr{छटा} (chaṭā) — a shade; a tint}

\label{lesson:MR-C535-chata}



\begin{quote}
Before the new one: say the Marathi for a sculpture, then the Marathi for a paintbrush.
\end{quote}

\subsection*{You'll want to know: \mr{छटा}}
\textbf{\mr{छटा}} — \emph{chaṭā} — \textquotedblleft{}a shade; a tint\textquotedblright{}.

\begin{quote}
\mr{आकाशात} \mr{निळ्या} \mr{रंगाच्या} \mr{खूप} \mr{छटा} \mr{होत्या}.
\end{quote}

\begin{quote}
\emph{ākāśāt niḷyā raṅgācyā khūp chaṭā hotyā.}
\end{quote}

\begin{quote}
There were many shades of blue in the sky.
\end{quote}

Feminine.

\begin{cousinweb}
\textbf{\mr{छटा}} (\emph{chaṭā}) is a tatsama, taken from Sanskrit \emph{chaṭā}, \textquotedblleft{}a mass, a lustre\textquotedblright{} (Monier-Williams).
\end{cousinweb}

\begin{grammarlens}[title={What you've built}]
One more word for describing what you lived and felt.
\end{grammarlens}

\subsection*{Your turn}
\begin{itemize}
\raggedright
  \item \emph{Say it:} \emph{chaṭā}
  \item \emph{Say it:} \emph{chaṭā}, once more
  \item \emph{Say it:} \emph{ākāśāt niḷyā raṅgācyā khūp chaṭā hotyā.} — There were many shades of blue in the sky.
  \item \emph{Recall:} say the Marathi for a sculpture, then the Marathi for a paintbrush, then say \emph{chaṭā} again
\end{itemize}

\begin{tcolorbox}[breakable,colback=teal!4,colframe=teal!35!black,title={Before you move on}]
What does \mr{छटा} mean? (\textquotedblleft{}A shade; a tint\textquotedblright{}.) Say it once more.
\end{tcolorbox}
\section[(dialogue)]{Looking back: from Hobbies and Leisure to Painting and Sculpture}

\label{lesson:MR-R535-review-528-535}



\begin{quote}
Say the Marathi for a hobby and a pastime.
\end{quote}

\subsection*{Your turn}
Say these aloud:

\begin{itemize}
\raggedright
  \item a hobby, a pastime, free time, entertainment, to be absorbed
  \item a match, a team captain, an umpire, a playing field, to overcome
  \item a medal, a trophy cup, a sporting best, a contestant, to carry off
  \item a batsman, a bowler, a runner, to hit away, to catch a ball
  \item chess, a game piece, a die for games, a crossword, hide and seek
  \item a workout, a yoga posture, a whistle, to bend, to leap
  \item mountaineering, swimming as a sport, a camp, a tent, a campfire
  \item painting as an art, an artist, a sculpture, a paintbrush, a shade
\end{itemize}

\begin{tcolorbox}[breakable,colback=teal!4,colframe=teal!35!black,title={Before you move on}]
A hobby? (\textbf{\mr{छंद}}, \emph{chand}.) A pastime? (\textbf{\mr{विरंगुळा}}, \emph{viraṅguḷā}.) Free time? (\textbf{\mr{फुरसत}}, \emph{phursat}.) Entertainment? (\textbf{\mr{करमणूक}}, \emph{karamṇūk}.) To be absorbed? (\textbf{\mr{दंग} \mr{होणे}}, \emph{daṅg hoṇe}.) A match? (\textbf{\mr{सामना}}, \emph{sāmnā}.) A team captain? (\textbf{\mr{कर्णधार}}, \emph{karṇadhār}.) An umpire? (\textbf{\mr{पंच}}, \emph{pañc}.) A playing field? (\textbf{\mr{मैदान}}, \emph{maidān}.) To overcome? (\textbf{\mr{मात} \mr{करणे}}, \emph{māt karṇe}.) A medal? (\textbf{\mr{पदक}}, \emph{padak}.) A trophy cup? (\textbf{\mr{चषक}}, \emph{caṣak}.) A sporting best? (\textbf{\mr{विक्रम}}, \emph{vikram}.) A contestant? (\textbf{\mr{स्पर्धक}}, \emph{spardhak}.) To carry off? (\textbf{\mr{पटकावणे}}, \emph{paṭkāvṇe}.) A batsman? (\textbf{\mr{फलंदाज}}, \emph{phalandāj}.) A bowler? (\textbf{\mr{गोलंदाज}}, \emph{golandāj}.) A runner? (\textbf{\mr{धावपटू}}, \emph{dhāvpaṭū}.) To hit away? (\textbf{\mr{टोलवणे}}, \emph{ṭolavṇe}.) To catch a ball? (\textbf{\mr{झेलणे}}, \emph{jhelṇe}.) Chess? (\textbf{\mr{बुद्धिबळ}}, \emph{buddhibaḷ}.) A game piece? (\textbf{\mr{सोंगटी}}, \emph{soṅgṭī}.) A die for games? (\textbf{\mr{फासा}}, \emph{phāsā}.) A crossword? (\textbf{\mr{शब्दकोडे}}, \emph{śabdakoḍe}.) Hide and seek? (\textbf{\mr{लपाछपी}}, \emph{lapāchpī}.) A workout? (\textbf{\mr{कसरत}}, \emph{kasrat}.) A yoga posture? (\textbf{\mr{योगासन}}, \emph{yogāsan}.) A whistle? (\textbf{\mr{शिट्टी}}, \emph{śiṭṭī}.) To bend? (\textbf{\mr{वाकणे}}, \emph{vākṇe}.) To leap? (\textbf{\mr{झेपावणे}}, \emph{jhepāvṇe}.) Mountaineering? (\textbf{\mr{गिर्यारोहण}}, \emph{giryārohaṇ}.) Swimming as a sport? (\textbf{\mr{जलतरण}}, \emph{jaltaraṇ}.) A camp? (\textbf{\mr{शिबिर}}, \emph{śibir}.) A tent? (\textbf{\mr{तंबू}}, \emph{tambū}.) A campfire? (\textbf{\mr{शेकोटी}}, \emph{śekoṭī}.) Painting as an art? (\textbf{\mr{चित्रकला}}, \emph{citrakalā}.) An artist? (\textbf{\mr{कलाकार}}, \emph{kalākār}.) A sculpture? (\textbf{\mr{शिल्प}}, \emph{śilpa}.) A paintbrush? (\textbf{\mr{कुंचला}}, \emph{kuñclā}.) A shade? (\textbf{\mr{छटा}}, \emph{chaṭā}.)
\end{tcolorbox}
